\documentclass[oneside]{sysuthesis}

\sysusetup[option]
  {
    % anonymous,
    type = bachelor
   % type = master
    % type = doctor
  }
\sysusetup[style]
  {
    % cjk-font  = windows,
    % math-font = newcm
  }
\sysusetup[info]
  {
    title        = {基于高斯过程仿真器核函数选择的\\ 宇宙学参数约束研究},
    title*       = {Cosmological Parameter Constraints from\\
                    Gaussian Process Emulators with Different Kernel Functions},
    keywords     = {宇宙学仿真器, 高斯过程, 核函数, 参数约束, MCMC, CSST},
    keywords*    = {cosmological emulator, Gaussian process, kernel function,
                    parameter constraint, MCMC, CSST},
    student-id   = {22344031},
    author       = {陈伟亮},
    author*      = {XX},
    department   = {物理与天文学院},
    department*  = {XX School},
    major        = {物理专业},
    major*       = {XX Major},
    supervisors  = {李霄栋~教授},
    supervisors* = {Prof.~XX},
    date         = {2026/04/10}
  }


% 参考文献设置: 使用 BibTeX
\sysusetup[bib]
  {
    backend  = bibtex,
    style    = gbt7714-numerical,
    resource = reference/sysuthesis-sample
  }

% 载入其他常用宏包
\usepackage{siunitx}        % 物理单位
\usepackage{subcaption}     % 图片并排
\usepackage{tabularray}     % 表格
\UseTblrLibrary{booktabs}   % 三线表样式


% ========== 【长表格全局设置区】 ==========
\DefTblrTemplate{contfoot-text}{default}{} % 强制清空跨页底部的英文提示
\DefTblrTemplate{conthead-text}{default}{} % 强制清空跨页顶部的英文提示
\SetTblrStyle{caption-tag}{font=\small}    % 强制长表格的“表号”（如：表4.1）为小字号且不加粗
\SetTblrStyle{caption-text}{font=\small}   % 强制长表格的“表名”为小字号且不加粗
% ==============================================
% ==============================================

\usepackage{float}          % 浮动体
% amsmath 已由 cls 加载；amssymb 和 bm 与 unicode-math 冲突，已由其替代
% 如需粗体数学符号，使用 \symbf{} 替代 \bm{}
% booktabs 已由 tabularray 的 booktabs 库提供，无需重复加载

% ========== 【图表公式编号：点号改短横线，符合中大格式要求】 ==========
\renewcommand{\thefigure}{\thechapter-\arabic{figure}}
\renewcommand{\thetable}{\thechapter-\arabic{table}}
\renewcommand{\theequation}{\thechapter-\arabic{equation}}
% =====================================================================



\begin{document}

\maketitle
\frontmatter

% !TeX root = ../thesis.tex
\begin{abstract}
  随着宇宙学大型巡天项目（如 CSST、DESI、Euclid 等）的推进，观测数据的精度和规模迅速提升。
  传统的两点关联函数和功率谱分析方法面临非线性尺度信息提取不足的挑战，亟需开发高精度的理论
  预测工具。宇宙学仿真器（emulator）通过在高精度 N 体模拟上训练机器学习模型，建立统计量与
  宇宙学参数之间的映射关系，以极低的计算成本实现对任意参数组合下统计量的快速预测，已成为
  该领域的主流方法。

  本文系统研究了高斯过程（GP）仿真器中核函数选择对预测精度和参数约束能力的影响。基于 129
  组覆盖 9 维宇宙学参数空间的高精度 N 体模拟数据，分别针对各向同性两点关联函数 $\xi(s)$ 和
  距离积分关联函数 $\xi(\mu)$（定义为将二维关联函数 $\xi(s,\mu)$ 对分离距离 $s$
  在 $0$ 至 $120~\mathrm{Mpc}/h$ 范围内积分，保留角度 $\mu$ 为自变量）构建仿真器，
  系统比较了六种核函数的性能：RBF 核、Mat\'ern$(\nu=5/2)$ 核、Mat\'ern$(\nu=3/2)$ 核、
  Mat\'ern$(\nu=1/2)$ 核、Rational Quadratic（RQ）核及 Mat\'ern$\times$RBF 乘积组合核。
  实验采用严格的公平对比策略：随机训练/测试划分、固定随机种子、相同超参数搜索预算，并通过
  留一交叉验证（LOO-CV）进行精度评估。

  结果表明：在 $\xi(s)$ 上，Mat\'ern 族核（$\nu=5/2$ 和 $3/2$）和 RQ 核均显著优于 RBF 核
  （配对 Wilcoxon 检验 $p<0.0001$，经 Bonferroni 校正后仍显著），而 Mat\'ern$(\nu=1/2)$
  核因过度粗糙表现最差（MAE${}=0.00821$，比 RBF 高 $45\%$）。Mat\'ern$(\nu=5/2)$ 与
  Mat\'ern$(\nu=3/2)$ 性能无统计显著差异（$p=0.81$），表明 $\nu=5/2$ 和 $3/2$ 均为合适
  选择。$\xi(\mu)$ 上核函数差异在无泄露条件下不显著。Mat\'ern$(\nu=5/2)$ 核的联合约束品质
  因子 ${\rm FoM}=58.9$（bootstrap 标准差 $0.2$），相比单独使用 $\xi(s)$ 提升 $25\%$。
  六种核的联合 FoM 从高到低依次为：Mat\'ern$(5/2)$ ($58.9$) ${}>{}$ Mat\'ern$(3/2)$
  ($58.3$) ${}>{}$ RQ ($57.7$) ${}>{}$ Mat\'ern$(1/2)$ ($56.3$) ${}>{}$ Combined
  ($50.4$) ${}>{}$ RBF ($49.3$)。实验通过 3-seed 重复性验证（MAE 种子间标准差
  ${<}0.0002$）和 Gelman--Rubin 收敛诊断（$\hat R<1.01$）确认了结果稳健性。本研究推荐
  Mat\'ern$(\nu=5/2)$ 核作为面向 CSST 宇宙学仿真器的默认选择。
\end{abstract}

\begin{abstract*}
  With the rapid advancement of large-scale cosmological surveys such as CSST, DESI,
  and Euclid, traditional two-point correlation function analyses face challenges in
  extracting nonlinear-scale information. Cosmological emulators, which use machine
  learning models trained on N-body simulations to map statistical observables to
  cosmological parameters, have become the mainstream approach.

  This thesis systematically investigates the impact of kernel function selection in
  Gaussian Process (GP) emulators on prediction accuracy and parameter constraint
  capability. Based on $129$ high-precision N-body simulations spanning a $9$-dimensional
  cosmological parameter space, emulators are constructed for both the isotropic
  two-point correlation function $\xi(s)$ and the distance-integrated correlation
  function $\xi(\mu)=\int_0^{120}\xi(s,\mu)\,ds$ (integrated over separations
  $0$--$120~\mathrm{Mpc}/h$). Six kernel functions are compared under a rigorous
  fair-comparison protocol: RBF, Mat\'ern$(\nu=5/2)$, Mat\'ern$(\nu=3/2)$,
  Mat\'ern$(\nu=1/2)$, Rational Quadratic, and Mat\'ern$\times$RBF product kernel,
  with randomized train/test splitting, independent random seeds, equal hyperparameter
  budgets via Optuna, and leave-one-out cross-validation (LOO-CV) for accuracy
  assessment.

  Results show: (1) For $\xi(s)$, the Mat\'ern family ($\nu=5/2$ and $3/2$) achieves
  LOO MAE ${\sim}0.0051$, significantly outperforming RBF ($0.00565$, $p<0.0001$) and
  RQ ($0.00539$, $p<0.001$); Mat\'ern$(\nu=1/2)$ performs worst (MAE ${}=0.00821$,
  too rough). Mat\'ern$(5/2)$ and Mat\'ern$(3/2)$ show no significant difference
  ($p=0.81$). (2) For $\xi(\mu)$, kernel differences are not robust under
  leakage-free conditions. (3) Joint FoM ranking: Mat\'ern$(5/2)$ ($58.9$) ${}>{}$
  Mat\'ern$(3/2)$ ($58.3$) ${}>{}$ RQ ($57.7$) ${}>{}$ Mat\'ern$(1/2)$ ($56.3$)
  ${}>{}$ Combined ($50.4$) ${}>{}$ RBF ($49.3$). (4) Results are highly reproducible
  across $3$ independent random seeds. (5) All MCMC chains pass Gelman--Rubin
  convergence diagnostics ($\hat R<1.01$). This work recommends Mat\'ern$(\nu=5/2)$
  as the default kernel for CSST-oriented cosmological emulators.
\end{abstract*}



\tableofcontents
\listoffigures
\listoftables

\mainmatter

\chapter{绪论}\label{chap:intro}

\section{研究背景与意义}\label{sec:bg}

宇宙学的核心任务之一是精确测量描述宇宙组成和演化的基本参数，包括冷暗物质密度参数
$\Omega_c$、物质密度扰动振幅 $\sigma_8$、暗能量状态方程参数 $w$ 等。21 世纪以来，
以宇宙微波背景辐射（CMB）、重子声学振荡（BAO）和弱引力透镜为代表的观测手段极大地推进了
精确宇宙学的发展\cite{Planck2020,Riess1998}。当前，以中国空间站望远镜（CSST）\cite{Zhan2011}、
暗能量光谱仪（DESI）、Euclid 空间望远镜为代表的新一代大型巡天项目即将或已经开始运行，
预计将覆盖数十亿星系，在前所未有的统计精度下探测宇宙大尺度结构。

在大尺度结构研究中，两点关联函数 $\xi(r)$ 及其傅里叶对偶——功率谱 $P(k)$
是描述物质密度场统计性质的基本工具。在线性理论框架下，这些统计量可以通过 Boltzmann 方程组
（如 CAMB、CLASS）精确计算。然而，在非线性尺度（$k>0.1~h/\mathrm{Mpc}$），微扰论失效，
必须依赖 N 体数值模拟来获取准确的理论预测。单次高分辨率 N 体模拟可能需要数千 CPU 核时，
在宇宙学参数空间中密集采样极其昂贵。

宇宙学仿真器（cosmological emulator）应运而生：在参数空间中通过拉丁超立方采样\cite{McKay1979}
选取有限组代表性样本，在每个样本点运行高精度 N 体模拟，再训练机器学习模型（通常是高斯过程）
插值预测任意参数下的统计量。代表性工作包括 Cosmic Emulator\cite{Heitmann2009}、
EuclidEmulator2\cite{Knabenhans2019} 等。

\section{国内外研究现状}\label{sec:related}

\subsection{宇宙学仿真器的发展}

高斯过程（GP）因其非参数特性和能提供预测不确定性估计的优势，是宇宙学仿真器的主流建模方法。
Heitmann 等人\cite{Heitmann2009}开发的 Cosmic Emulator 基于 37 组 N 体模拟构建了物质功率谱
的 GP 仿真器，开创了该领域。Knabenhans 等人\cite{Knabenhans2019}在 EuclidEmulator2 中
进一步将仿真器推广到含中微子质量的高维参数空间，精度达 $1\%$ 以内。Nishimichi 等
人\cite{Nishimichi2019}基于 Dark Quest 项目构建了暗物质晕统计量的仿真器，覆盖了非线性结构
形成的多个方面。Villaescusa-Navarro 等人\cite{Villaescusa2021}发起的 CAMELS 项目提供了
数千组流体动力学模拟，极大地推动了仿真器方法论的发展。

近年来，基于神经网络的仿真器方法快速发展。Arico 等人\cite{Arico2021}使用神经网络构建了
Bacco 仿真器，在功率谱预测上达到亚百分比精度。Spurio Mancini 等人\cite{Mancini2022}开发的
CosmoPower 采用深度学习直接学习功率谱与参数的映射，计算速度比传统 Boltzmann 求解器快
数个数量级。然而，GP 仿真器在小样本（${<}1000$ 组模拟）场景下仍具优势，且能自然提供预测
不确定性。

\subsection{GP 核函数选择的研究}

在机器学习领域，核函数选择是 GP 回归的核心问题之一。Rasmussen 和 Williams\cite{Rasmussen2006}
的经典著作系统介绍了 GP 回归和主要核函数族的性质。Stein\cite{Stein1999}从空间统计学角度
论证了 Mat\'ern 核族相比 RBF 核在实际数据建模中的优势。Duvenaud\cite{Duvenaud2014}探索了
核函数的自动构建方法，通过组合基础核函数搜索最优核结构。然而，这些通用研究未针对宇宙学
数据的特殊性（高维参数空间、弱信号、强退化）进行考察。

\subsection{面向 CSST 的仿真器研究}

中国空间站望远镜（CSST）\cite{Zhan2011}是我国旗舰级空间天文项目，预计覆盖约 17500 平方度
天区。Gong 等人\cite{Gong2019}和 Lin 等人\cite{Lin2022}围绕 CSST 的弱引力透镜和星系成团性
开展了预测研究。Zhang 等人\cite{Zhang2023}基于 129 组高精度 N 体模拟构建了面向 CSST 的两点
关联函数仿真器，本文工作即在此基础上开展。然而，包括 Zhang 等人在内的现有工作大多采用固定
的 RBF 核函数，对核函数选择缺乏系统性的比较研究，尤其缺乏统计显著性检验层面的定量分析。
不同核函数隐含了对目标函数光滑性的不同先验假设，这直接影响仿真器的预测精度和参数约束能力。
本文旨在填补这一空白，为 CSST 数据分析管线中的仿真器设计提供定量依据。

\section{研究内容与创新点}\label{sec:contrib}

本文主要研究内容：（1）基于 129 组高精度 N 体模拟，构建面向 CSST 的 GP 宇宙学仿真器；
（2）系统比较六种核函数：RBF、Mat\'ern$(\nu=5/2)$、Mat\'ern$(\nu=3/2)$、
Mat\'ern$(\nu=1/2)$、Rational Quadratic（RQ）和 Mat\'ern$\times$RBF 乘积组合核；
（3）通过 MCMC 采样获得宇宙学参数约束，评估核函数选择的影响；（4）进行全面的统计诊断分析。

创新点：（1）首次对宇宙学仿真器核函数进行含统计显著性检验的系统比较；（2）严格公平对比策略
（随机划分、独立种子、等预算、3-seed 重复性验证）；（3）从仿真器精度、参数约束、MCMC 收敛
等多维度全面评估；（4）提供可操作的核函数选择建议，为 CSST 数据分析提供技术支撑。

本文的总体技术路线如 \autoref{fig:roadmap} 所示：从 N 体模拟数据出发，经预处理后训练
六种核函数的 GP 仿真器，通过 Optuna 超参数优化、LOO-CV 泛化评估和 Wilcoxon 统计检验
进行公平的模型选择，最终用 MCMC 采样获得宇宙学参数约束并计算 FoM 品质因子。

\begin{figure}[H]
  \centering
  \includegraphics[width=0.95\linewidth]{01_技术路线图.png}
  \caption{总体技术路线图}
  \label{fig:roadmap}
\end{figure}

\section{论文结构}\label{sec:outline}

本文共六章：第一章绪论；第二章宇宙学理论基础（含 FLRW 度规、结构增长、RSD）；第三章研究方法
（GP 回归、核函数、公平对比策略、MCMC）；第四章实验结果（精度、显著性检验、MAE 物理含义、
MCMC 约束、收敛诊断、重复性验证）；第五章讨论；第六章总结与展望。

\chapter{宇宙学理论基础}\label{chap:theory}

\section{宇宙膨胀与 Friedmann 方程}\label{sec:friedmann}

在广义相对论框架下，均匀各向同性宇宙（宇宙学原理）的时空度规由 Friedmann--Lema\^{i}tre--
Robertson--Walker（FLRW）度规描述。宇宙膨胀演化由 Friedmann 方程控制。在低红移（$z<10$）下，
辐射能量密度可忽略，方程简化为：
\begin{equation}\label{eq:friedmann}
  H(a)^2 = H_0^2 \left[\, \Omega_m\, a^{-3} + \Omega_{\rm DE}(a) \,\right],
\end{equation}
其中 $H(a)=(da/dt)/a$ 是 Hubble 参数，$a$ 是标度因子（今天 $a=1$），$\Omega_m$ 是物质密度
参数，$\Omega_{\rm DE}(a)$ 是暗能量密度的标度因子依赖项。本文采用 CPL
参数化\cite{Chevallier2001,Linder2003}描述暗能量状态方程：$w(a)=w_0+w_a(1-a)$，对应的
暗能量密度演化为：
\begin{equation}\label{eq:de-density}
  \frac{\rho_{\rm DE}(a)}{\rho_{\rm DE}(1)}
    = a^{-3(1+w_0+w_a)}\, \exp\!\left[-3 w_a (1-a)\right].
\end{equation}
当 $w_a=0$ 时退化为恒定 $w$ 模型，当 $w_0=-1$ 且 $w_a=0$ 时退化为 $\Lambda$CDM 模型
（宇宙学常数）。完整 Friedmann 方程还应包含辐射项 $\Omega_r\, a^{-4}$，但在低红移下远小于
物质项，通常忽略。

\section{物质密度扰动与结构增长}\label{sec:growth}

宇宙大尺度结构起源于早期宇宙的微小密度扰动。定义密度对比
$\delta(\symbf{x})=[\rho(\symbf{x})-\bar\rho]/\bar\rho$，在线性微扰论框架下，其傅里叶模式
$\delta_{\symbf{k}}$ 满足：
\begin{equation}\label{eq:linear-growth}
  \frac{\mathrm d^2 \delta_{\symbf{k}}}{\mathrm d t^2}
    + 2H \,\frac{\mathrm d \delta_{\symbf{k}}}{\mathrm d t}
    - 4\pi G\, \bar\rho\, \delta_{\symbf{k}} = 0.
\end{equation}
该方程的增长解定义了线性增长因子 $D(a)$。参数 $\sigma_8$ 定义为在半径 $R=8~\mathrm{Mpc}/h$
球内的密度扰动均方根：
\begin{equation}\label{eq:sigma8}
  \sigma_8^2 = \frac{1}{2\pi^2}
    \int_0^{\infty} k^2\, P(k)\, W^2(kR)\, \mathrm d k,
\end{equation}
其中 $W(kR)$ 是半径为 $R=8~\mathrm{Mpc}/h$ 的顶帽窗函数，$P(k)$ 是物质功率谱。$\sigma_8$ 与
功率谱的归一化直接相关，是本文约束的核心参数之一。

\section{两点关联函数与功率谱}\label{sec:two-point}

两点关联函数 $\xi(r)$ 定义为密度场的二阶统计量：
\begin{equation}\label{eq:xi-def}
  \xi(r) = \langle \delta(\symbf{x})\, \delta(\symbf{x}+\symbf{r}) \rangle_{\symbf{x}},
\end{equation}
其中 $\langle\cdot\rangle_{\symbf{x}}$ 表示对位置 $\symbf{x}$ 的系综平均。$\xi(r)$ 与物质功率谱
$P(k)$ 通过傅里叶变换互为对偶：
\begin{equation}\label{eq:xi-fft}
  \xi(r) = \frac{1}{2\pi^2}
    \int_0^{\infty} P(k)\, k^2\, \frac{\sin(kr)}{kr}\, \mathrm d k.
\end{equation}

\section{红移空间畸变与 Kaiser 公式}\label{sec:rsd}

星系的本动速度（peculiar velocity）沿视线方向的分量会使红移测距产生系统偏差，
即红移空间畸变（Redshift Space Distortion，RSD）效应。在线性理论下，
Kaiser\cite{Kaiser1987}推导出红移空间功率谱与实空间功率谱的关系：
\begin{equation}\label{eq:kaiser}
  P^{s}(k,\mu_k) = \left(1+\beta\,\mu_k^2\right)^2 P^{r}(k),
\end{equation}
其中 $\mu_k=\cos(\theta_k)$ 是傅里叶空间中波矢 $\symbf{k}$ 与视线方向的夹角余弦，
$\beta=f/b$ 是 RSD 参数，$f\approx \Omega_m(z)^{0.55}$ 是线性增长率\cite{Linder2005}，
$b$ 是星系偏差参数。在配置空间中，红移空间的二维关联函数 $\xi(s,\mu)$ 依赖于分离向量
与视线方向的夹角 $\mu=\cos(\theta)$。需注意：Kaiser 公式中的 $\mu_k$ 是傅里叶空间角度
余弦，而 $\xi(s,\mu)$ 中的 $\mu$ 是配置空间角度余弦，两者物理含义不同但取值范围相同
（$[0,1]$），本文后续 $\xi(\mu)$ 均指配置空间中的 $\mu$。配置空间下的 Kaiser 公式
及其多极展开综述参见 Hamilton\cite{Hamilton1998}。本文所用的 $\xi(\mu)$ 定义为
将 $\xi(s,\mu)$ 在分离距离 $s$ 的 $0$ 到 $120~\mathrm{Mpc}/h$ 范围上积分：
\begin{equation}\label{eq:xi-mu}
  \xi(\mu) = \int_{0}^{120~\mathrm{Mpc}/h} \xi(s,\mu)\, \mathrm d s.
\end{equation}
共 97 个 $\mu$ bin，覆盖 $\mu\in[0,1]$，$\mu$ 间距约 $0.0104$。$\xi(\mu)$ 保留了 RSD 效应
对 $\mu$ 的依赖性，包含结构增长率 $f$ 的信息，与 $\xi(s)$ 提供互补的宇宙学约束。

\section{N 体数值模拟}\label{sec:nbody}

在非线性尺度（$k>0.1~h/\mathrm{Mpc}$）上必须依赖 N 体数值模拟。N 体模拟从初始功率谱出发，
通过 PM（Particle-Mesh）或 Tree-PM 等算法演化暗物质粒子，直接计算密度场的非线性演化。
本文使用的模拟数据来自 Zhang 等人\cite{Zhang2023}的工作：采用拉丁超立方采样\cite{McKay1979}
在 9 维参数空间（$\Omega_c$、$\sigma_8$、$H_0$、$\Omega_b$、$n_s$、$A_s$、$w$、$w_a$、
$m_\nu$）中生成 129 组采样点，每组使用 $1024^3$ 个粒子在 $1000~\mathrm{Mpc}/h$ 的模拟箱
中演化，统计量分别在 $s=0.5$ 到 $119.5~\mathrm{Mpc}/h$（$\xi(s)$，120 个 bin）和
$\mu=0$ 到 $1$（$\xi(\mu)$，97 个 bin）上计算。fiducial 参数为：$\Omega_c=0.26072$、
$\sigma_8=0.81006$、$H_0=67.66~\mathrm{km/s/Mpc}$、$\Omega_b=0.04898$、$n_s=0.9665$、
$A_s=2.105\times 10^{-9}$、$w=-1.0$、$w_a=0.0$、$m_\nu=0.06~\mathrm{eV}$。
协方差矩阵从边长 $400~\mathrm{Mpc}/h$ 的模拟箱中估计。需要说明的是，MCMC 约束中使用的
``观测值''（ground truth）来自边长 $600~\mathrm{Mpc}/h$ 的独立模拟箱
（${\rm gt\_boxsize}=600$），与训练数据（$1000~\mathrm{Mpc}/h$ 箱）不同。由于宇宙方差和
分辨率差异，同一参数下不同箱体的关联函数存在系统偏差，约 $51.5\%$ 的 $\xi(\mu)$ bin 和
$16.7\%$ 的 $\xi(s)$ bin 落在训练数据的取值范围之外。这意味着仿真器在这些 bin 上需要外推
而非内插，MCMC 约束的绝对结果受此影响（详见 \autoref{sec:omega-c-bias} 讨论），但不影响
不同核函数间的相对性能比较。此外，上游数据管线对每个模拟的关联函数做了逐样本均值归一化
（$\xi\to\xi/\mathrm{mean}(\xi)$），使得每个样本在所有 bin 上的均值恰好为 $1$。这意味着
仿真器学习的是关联函数的形状而非绝对幅度。由于 $\sigma_8$ 主要控制关联函数的整体振幅，
归一化后 $\sigma_8$ 的信号被削弱，这是 $\xi(\mu)$ 约束能力较弱和 $\sigma_8$ 偏移较大的
深层原因之一。

\chapter{研究方法}\label{chap:method}

\section{高斯过程回归}\label{sec:gp}

高斯过程（GP）是一种非参数贝叶斯回归方法，将函数空间上的先验定义为
$f(\symbf{x})\sim\mathcal{GP}(m(\symbf{x}),\,k(\symbf{x},\symbf{x}'))$，其中 $m(\symbf{x})$ 为均值函数
（取为零），$k(\symbf{x},\symbf{x}')$ 为核函数。给定训练数据 $\{X,\symbf{y}\}$，GP 的预测分布为
解析高斯分布：
\begin{align}
  f^{*}(\symbf{x}^{*}) &\sim \mathcal N(\mu^{*},\, {\sigma^{*}}^2),
    \label{eq:gp-pred} \\
  \mu^{*} &= k(\symbf{x}^{*},X)\, \left[K + \sigma_n^2 I\right]^{-1} \symbf{y},
    \label{eq:gp-mean} \\
  {\sigma^{*}}^2 &= k(\symbf{x}^{*},\symbf{x}^{*})
    - k(\symbf{x}^{*},X)\, \left[K + \sigma_n^2 I\right]^{-1} k(X,\symbf{x}^{*}),
    \label{eq:gp-var}
\end{align}
其中 $K$ 是训练数据的核矩阵，$\sigma_n^2$ 是噪声方差。GP 回归的计算复杂度为 $O(n^3)$
（核矩阵求逆），对本文的 129 个训练样本完全可行（单次训练约 $2$ 秒）。本文采用
scikit-learn\cite{Pedregosa2011}的多输出 GP 架构：单个 GP 接收 9 维参数输入，同时预测
所有输出 bin（$\xi(s)$ 为 120 维，$\xi(\mu)$ 为 97 维），所有 bin 共享核函数超参数和
核矩阵 $K$，每个 bin 通过独立的回归权重 $\alpha_j=K^{-1}\symbf{y}_j$ 分别计算预测值。
这不同于逐 bin 独立训练 120/97 个 GP 的方式。对于更大规模（1000+ 样本）的场景，可考虑
稀疏 GP（FITC/SVGP）或随机傅里叶特征等近似方法。

GP 仿真器的训练和预测流程原理见 \autoref{fig:gp-principle}：训练阶段通过核函数计算
协方差矩阵并拟合回归模型，预测阶段利用训练好的模型对新参数给出均值预测和不确定度估计。

\begin{figure}[H]
  \centering
  \includegraphics[width=0.95\linewidth]{02_gp_principle.png}
  \caption{高斯过程仿真器原理示意图}
  \label{fig:gp-principle}
\end{figure}

\section{核函数}\label{sec:kernels}

本文系统比较以下六种核函数：

\paragraph{（1）RBF 核（平方指数核）}
\begin{equation}\label{eq:k-rbf}
  k_{\rm RBF}(\symbf{x},\symbf{x}') = \sigma^2\, \exp\!\left(
    -\frac{\lVert \symbf{x}-\symbf{x}' \rVert^2}{2\ell^{2}} \right).
\end{equation}
RBF 核生成无限可微的极光滑函数，是 GP 回归最常用的核函数。

\paragraph{（2）Mat\'ern$(\nu=5/2)$ 核}
\begin{equation}\label{eq:k-m52}
  k_{\rm M52}(\symbf{x},\symbf{x}') = \sigma^2 \left(
    1 + \frac{\sqrt{5}\,r}{\ell} + \frac{5 r^2}{3\ell^2}
  \right) \exp\!\left(-\frac{\sqrt{5}\,r}{\ell}\right),
  \quad r=\lVert \symbf{x}-\symbf{x}' \rVert.
\end{equation}
生成二阶可微函数，在光滑性与灵活性间取得良好平衡，是 GP 文献中最常用的 Mat\'ern
变体\cite{Rasmussen2006}。

\paragraph{（3）Mat\'ern$(\nu=3/2)$ 核}
\begin{equation}\label{eq:k-m32}
  k_{\rm M32}(\symbf{x},\symbf{x}') = \sigma^2 \left(1 + \frac{\sqrt{3}\,r}{\ell}\right)
    \exp\!\left(-\frac{\sqrt{3}\,r}{\ell}\right).
\end{equation}
生成一阶可微函数，比 $\nu=5/2$ 更灵活但也更``粗糙''。

\paragraph{（4）Mat\'ern$(\nu=1/2)$ 核（指数核）}
\begin{equation}\label{eq:k-m12}
  k_{\rm M12}(\symbf{x},\symbf{x}') = \sigma^2\, \exp\!\left(-\frac{r}{\ell}\right).
\end{equation}
生成连续但不可微函数，是最粗糙的 Mat\'ern 变体。$\nu=1/2$、$3/2$、$5/2$、$\infty$（RBF）
构成光滑度从低到高的完整序列，本文通过同时测试这四种核来系统考察光滑度对仿真器性能的影响。

\paragraph{（5）Rational Quadratic（RQ）核}
\begin{equation}\label{eq:k-rq}
  k_{\rm RQ}(\symbf{x},\symbf{x}') = \sigma^2 \left(
    1 + \frac{r^2}{2\,\alpha\,\ell^2}
  \right)^{-\alpha}.
\end{equation}
RQ 核等价于不同长度尺度 RBF 核的无穷混合，额外参数 $\alpha$ 控制混合权重分布。当
$\alpha\to\infty$ 时退化为 RBF 核。RQ 核可建模多尺度特征而不需要显式的乘积组合。

\paragraph{（6）乘积组合核}
\begin{equation}\label{eq:k-comb}
  k_{\rm comb}(\symbf{x},\symbf{x}') = k_{\rm M52}(\symbf{x},\symbf{x}')\cdot k_{\rm RBF}(\symbf{x},\symbf{x}').
\end{equation}
将两种核相乘可建模多个特征长度尺度（multi-length-scale），理论上表达力更强，但引入额外
超参数，在有限训练样本下可能导致优化困难。

所有核均附加 WhiteKernel 项 $k_{\rm white}=\sigma_n^2\,\delta(\symbf{x},\symbf{x}')$ 建模观测噪声，
完整核形式为 $k_{\rm total}=C\cdot k_{\rm main}+k_{\rm white}$，$C$ 为常数核。

上述六种核函数在光滑度谱上构成清晰的序列（\autoref{fig:kernel-smoothness}）：Mat\'ern
族核函数通过参数 $\nu$ 控制光滑度，$\nu=1/2$（最粗糙）、$3/2$、$5/2$ 到 RBF（无穷光滑）。
本文的发现之一是光滑度与仿真器性能呈非单调关系——过低和过高的光滑度均不利于宇宙学
仿真，中等光滑度的 Mat\'ern$(5/2)$ 和 Mat\'ern$(3/2)$ 表现最优。

\begin{figure}[H]
  \centering
  \includegraphics[width=0.95\linewidth]{03_kernel_smoothness.png}
  \caption{核函数光滑度与仿真器性能的非单调关系示意图}
  \label{fig:kernel-smoothness}
\end{figure}

\section{超参数优化}\label{sec:hpo}

核函数的超参数优化分为两层：内层由 sklearn 通过最大化边际似然自动优化所有核超参数
（$\ell$、信号方差、噪声水平等），从多个随机起点重启（${\rm n\_restarts}=10$）；
外层由 Optuna 搜索核参数初始值和正则化参数 $\alpha$（唯一不被内层优化覆盖的参数），
以测试集 MAE 为目标。本文采用 Optuna\cite{Akiba2019}自动超参数搜索框架，使用树结构
Parzen 估计（TPE）采样策略，在以下空间搜索：常数核振幅 $C\in[10^{-3},10^{4}]$
（对数均匀）、白噪声水平 $\sigma_n^2\in[10^{-6},10^{-2}]$（对数均匀）、正则化参数
$\alpha\in[10^{-6},10^{-2}]$（对数均匀）。Optuna 优化以测试集预测 MAE 为目标，搜索最优
超参数配置；完成后，以最优超参数重新训练模型，再在全部 129 个样本上运行 LOO-CV
计算最终精度指标。

需要指出方法论上的两处数据泄露：（1）Optuna 直接以测试集 MAE 为优化目标，
这是测试集参与模型选择的直接泄露（非``间接''泄露），严格做法应采用嵌套交叉验证
（nested CV），在训练集内部用验证集选超参数。（2）数据归一化的均值和标准差在全部 129 个
样本上计算（包含测试集），之后才进行训练/测试划分，测试集的分布信息泄露到了预处理阶段。
本文未采用嵌套 CV 主要因为计算成本考量（$129\times 20=2580$ 次 GP 训练）。但由于 LOO-CV
不重新优化超参数（${\rm optimizer}={\rm None}$），3-seed 验证显示超参数优化已完全收敛
（$\xi(s)$ MAE 种子间标准差 ${<}10^{-6}$）。关于预处理泄露：全集归一化使测试集的分布信息
（均值和标准差）部分泄露到训练阶段，其影响程度取决于测试集与训练集分布的差异大小
（而非简单地与样本比例 $1/129$ 成正比）。无泄露对照实验直接量化了这一影响。

为量化泄露影响，本文进行了无泄露对照实验（归一化仅在训练集 fit，Optuna 用训练集内部
验证集）。结果表明：（1）$\xi(s)$ 的 LOO MAE 升高约 $10\%$（从 $0.00514$ 升至 $0.00593$，
Mat\'ern 核），但核函数排序不变（Mat\'ern 仍优于 RBF），且独立测试集 MAE（$0.00551$）与
有泄露版本接近，说明泄露虽使 LOO MAE 偏乐观，但核函数性能排序结论稳健。
（2）$\xi(\mu)$ 的三种核 MAE 在无泄露条件下几乎相同（$0.01530$--$0.01551$），Mat\'ern
不再显著优于 RBF，提示 $\xi(\mu)$ 上的核函数差异结论完全受泄露影响。因此，本文
的核函数推荐结论仅基于 $\xi(s)$ 上的 MAE 比较，$\xi(\mu)$ 的核函数排序（含显著性检验和
联合 FoM 的一部分）不作为推荐依据。每种核函数固定 20 次 Optuna 试验
（${\rm n\_trials}=20$）。

双层超参数优化的完整流程如 \autoref{fig:hpo-flow} 所示：外层 Optuna TPE 采样器在 20 次
试验中智能搜索 ${\rm constant\_value}$、${\rm noise\_level}$ 和 $\alpha$ 三个超参数；
每次试验内部，sklearn 的 L-BFGS-B 优化器从 10 个随机起点出发，最大化边际对数似然以优化
$\ell$ 等核参数；最终以测试集 MAE 作为 Optuna 的目标值，反馈至 TPE 模型指导下一轮搜索。

\begin{figure}[H]
  \centering
  \includegraphics[width=0.95\linewidth]{04_hpo_flow.png}
  \caption{双层超参数优化流程图（外层 Optuna TPE + 内层 sklearn L-BFGS-B）}
  \label{fig:hpo-flow}
\end{figure}

\section{公平对比策略}\label{sec:fairness}

\textbf{（1）随机训练/测试划分。} 固定随机种子将 129 组样本随机打乱，按 $80{:}20$ 划分为
训练集（104 组）和测试集（25 组），避免按序截取引入的偏差（原始数据第 1 行为 fiducial
参数点）。

\textbf{（2）控制随机性。} 为避免搜索路径的随机性干扰核函数比较，六种核分别使用固定的
Optuna 随机种子（RBF: $42$，Mat\'ern$(5/2)$: $142$，Combined: $242$，Mat\'ern$(3/2)$: $342$，
Mat\'ern$(1/2)$: $442$，RQ: $542$），MCMC 采样使用相同的初始 walker 位置
（${\rm np.random.seed}(42)$）。\autoref{sec:reproducibility} 的 3-seed 稳定性测试验证了
结果对种子选择的不敏感性（$\xi(s)$ MAE 种子间标准差 ${<}10^{-6}$），证明种子差异不影响结论。

\textbf{（3）等量搜索预算。} 所有核函数使用相同的 Optuna 试验次数（${\rm n\_trials}=20$）
和 GP 优化器重启次数（${\rm n\_restarts}=10$），保证计算资源分配均等。

\textbf{（4）统一数据预处理。} 所有统计量使用相同的归一化策略（减均值除标准差），确保 GP
对不同统计量的输入/输出数据尺度一致。这一步骤对 $\xi(\mu)$ 尤为关键：$\xi(\mu)$ 的原始
变化范围仅为 $0.87$--$1.28$，归一化后 GP 才能有效区分不同核函数的性能。

\section{精度评估：留一交叉验证}\label{sec:loo}

本文采用留一交叉验证（LOO-CV）作为主要精度指标。LOO-CV 在全部 129 个样本（包含训练集和
测试集）上进行：依次留出每个样本，用剩余 128 个样本训练 GP（使用 Optuna 已优化好的
超参数，不重新搜索），预测留出样本，计算预测误差。需要说明：LOO-CV 评估范围包含了 Optuna
搜索阶段使用过的测试集样本，这与 \autoref{sec:hpo} 讨论的数据泄露有关。其优势在于：
（1）不浪费数据，全部 129 个样本均参与评估；（2）对小样本量（129 组）提供稳定的泛化性能
估计。
\begin{equation}\label{eq:loo-mae}
  \mathrm{LOO~MAE}
    = \frac{1}{N} \sum_{i=1}^{N} \frac{1}{n_{\rm bin}}
       \sum_{j=1}^{n_{\rm bin}}
       \left| y_{ij}^{\rm pred} - y_{ij}^{\rm true} \right|,
\end{equation}
同时报告每个样本 MAE 的标准误差（${\rm SE}=\mathrm{std}/\sqrt N$）。

\section{MCMC 参数约束}\label{sec:mcmc}

使用仿射不变采样器 emcee\cite{ForemanMackey2013}进行 MCMC 参数约束。每一步提议参数
$(\Omega_c,\sigma_8)$，调用仿真器预测统计量，与``观测值''（fiducial 模拟结果）通过高斯
似然比较：
\begin{equation}\label{eq:lnL}
  \ln \mathcal{L} = -\tfrac{1}{2}\,
    \left[\symbf{d} - \symbf{m}(\symbf{\theta})\right]^{\!\top}
    C^{-1}
    \left[\symbf{d} - \symbf{m}(\symbf{\theta})\right],
\end{equation}
其中 $\symbf{d}$ 为观测数据向量，$\symbf{m}(\symbf{\theta})$ 为仿真器预测，$C$ 为协方差矩阵
（从边长 $400~\mathrm{Mpc}/h$ 的模拟箱估计）。

需要特别说明两点局限性：（1）本文仅拟合 $\Omega_c$ 和 $\sigma_8$ 两个参数，其余 7 个固定
在 fiducial 值。实际巡天分析需边缘化所有参数，FoM 值将显著降低；因此本文 FoM 值仅作为
不同核函数的相对性能比较基准，不代表 CSST 的实际预期约束能力。
（2）协方差矩阵从 $400~\mathrm{Mpc}/h$ 箱估计但仿真器数据来自 $1000~\mathrm{Mpc}/h$ 箱，
体积不匹配导致统计误差被高估（见 \autoref{sec:omega-c-bias} 节讨论）。

MCMC 配置：40 条 walker 链，5000 步，2000 步 burn-in，约束品质因子
\begin{equation}\label{eq:fom}
  {\rm FoM} = \frac{1}{\pi\, \sqrt{\det(\mathrm{Cov})}\cdot 2.3},
\end{equation}
其中 $2.3$ 对应二维参数空间中 $68.3\%$ 置信区间的 $\Delta\chi^2$ 阈值（$1\sigma$ 等值线）。
需要特别说明的是，由于协方差矩阵来自体积不匹配的模拟箱（见
\autoref{sec:omega-c-bias}），本文 FoM 的绝对数值不具有直接物理意义，不代表 CSST 实际
约束能力；FoM 仅用于六种核函数之间的相对排序比较，所有在不同核函数间共享同一协方差矩阵
的系统误差会被消去，不影响相对结论。

MCMC 流程如 \autoref{fig:mcmc-flow} 所示：40 条 Markov 链从初始参数出发，在每一步中依次
完成：提议新参数、先验范围检查、GP 仿真器毫秒级预测、似然函数计算和 Metropolis 接受/
拒绝判断。经过 5000 步迭代后，丢弃前 2000 步 burn-in 样本，对剩余样本进行
Gelman--Rubin\cite{Gelman1992}收敛诊断、后验分布提取、等高线图绘制和 FoM 计算。本文在
三种约束模式下分别运行 MCMC：仅 $\xi(s)$、仅 $\xi(\mu)$、以及联合 $\xi(s)+\xi(\mu)$。

\begin{figure}[H]
  \centering
  \includegraphics[width=0.95\linewidth]{05_mcmc_flow.png}
  \caption{MCMC 宇宙学参数约束流程图}
  \label{fig:mcmc-flow}
\end{figure}

\chapter{实验结果}\label{chap:results}

\section{参数空间采样}\label{sec:sampling}

\autoref{fig:sampling} 展示了 129 组 N 体模拟在 $\Omega_c$--$\sigma_8$ 平面上的拉丁超立方
采样分布。采样点均匀覆盖参数空间，红色星号标记 fiducial 参数位置（$\Omega_c=0.26072$，
$\sigma_8=0.81006$）。

\begin{figure}[H]
  \centering
  \includegraphics[width=0.85\linewidth]{image1.png}
  \caption{129 组 N 体模拟的参数空间采样分布（$\Omega_c$--$\sigma_8$ 投影）}
  \label{fig:sampling}
\end{figure}

\section{仿真器精度比较}\label{sec:accuracy}

\autoref{fig:loo-mae} 展示了六种核函数在 $\xi(s)$ 和 $\xi(\mu)$ 上的 LOO MAE 对比。
\autoref{tab:loo} 列出详细精度指标。

\begin{figure}[H]
  \centering
  \includegraphics[width=0.95\linewidth]{image2.png}
  \caption{六种核函数的 LOO MAE 对比}
  \label{fig:loo-mae}
\end{figure}

\begin{table}[H]
  \caption{六种核函数的 LOO-CV 精度（MAE $\pm$ SE）}
  \label{tab:loo}
  \centering
  \begin{tblr}{lcccc}
    \toprule
    核函数             & $\xi(s)$ MAE & $\xi(s)$ SE & $\xi(\mu)$ MAE & $\xi(\mu)$ SE \\
    \midrule
    RBF                & $0.00565$    & $0.00029$   & $0.01568$      & $0.00017$     \\
    Mat\'ern$(\nu=5/2)$ & $0.00514$    & $0.00025$   & $0.01532$      & $0.00015$     \\
    Mat\'ern$(\nu=3/2)$ & $0.00511$    & $0.00024$   & $0.01533$      & $0.00015$     \\
    Mat\'ern$(\nu=1/2)$ & $0.00821$    & $0.00052$   & $0.01534$      & $0.00015$     \\
    RQ                 & $0.00539$    & $0.00027$   & $0.01530$      & $0.00015$     \\
    Combined (Mat\'ern $\times$ RBF) & $0.00511$ & $0.00024$ & $0.01526$ & $0.00015$ \\
    \bottomrule
  \end{tblr}
\end{table}

\autoref{fig:optuna} 展示了 Optuna 超参数搜索过程中最优 MAE 随试验次数的收敛情况。
对于 $\xi(s)$，六种核均在前 $5$ 次试验内快速收敛，后续试验改善有限，说明 20 次试验已
充分。对于 $\xi(\mu)$，收敛同样迅速，且各核最终 MAE 差异较小。

\begin{figure}[H]
  \centering
  \includegraphics[width=0.95\linewidth]{image3.png}
  \caption{Optuna 超参数搜索收敛曲线（最优 MAE vs 试验次数）}
  \label{fig:optuna}
\end{figure}

\autoref{tab:best-kernels} 列出了 Optuna 搜索得到的最优核参数，便于读者复现和理解各核函数
的行为特征。

\begin{table}[H]
  \caption{各仿真器最优核参数（六种核函数）}
  \label{tab:best-kernels}
  \centering
  \footnotesize
  \begin{tblr}{llllc}
    \toprule
    核函数             & 统计量      & 信号核                                  & WhiteKernel 噪声 & $\alpha$    \\
    \midrule
    RBF                 & $\xi(s)$    & $3.88^{2}\cdot{\rm RBF}(\ell=11.5)$        & $0.00688$       & \num{1.27e-5} \\
    Mat\'ern$(5/2)$      & $\xi(s)$    & $5.53^{2}\cdot\mathrm{Mat\acute{e}rn}(\ell=21.4,\nu=5/2)$ & $0.00481$ & \num{4.97e-5} \\
    Mat\'ern$(3/2)$      & $\xi(s)$    & $11.4^{2}\cdot\mathrm{Mat\acute{e}rn}(\ell=60.5,\nu=3/2)$ & $0.00126$ & \num{8.39e-6} \\
    Mat\'ern$(1/2)$      & $\xi(s)$    & $3.63^{2}\cdot\mathrm{Mat\acute{e}rn}(\ell=281,\nu=1/2)$  & \num{1.0e-5} & \num{1.02e-6} \\
    RQ                  & $\xi(s)$    & $5.00^{2}\cdot{\rm RQ}(\alpha=1.06,\ell=15.2)$ & $0.00303$ & \num{3.26e-3} \\
    Combined            & $\xi(s)$    & $4.76^{2}\cdot\mathrm{Mat\acute{e}rn}(\ell=19.6)\cdot{\rm RBF}(\ell=\num{9.8e4})$ & $0.00360$ & \num{1.09e-6} \\
    \midrule
    RBF                 & $\xi(\mu)$  & $0.026^{2}\cdot{\rm RBF}(\ell=0.011)$       & $0.993$         & \num{8.47e-4} \\
    Mat\'ern$(5/2)$      & $\xi(\mu)$  & $0.425^{2}\cdot\mathrm{Mat\acute{e}rn}(\ell=2.21,\nu=5/2)$ & $0.832$ & \num{2.55e-4} \\
    Mat\'ern$(3/2)$      & $\xi(\mu)$  & $0.457^{2}\cdot\mathrm{Mat\acute{e}rn}(\ell=2.08,\nu=3/2)$ & $0.795$ & \num{8.44e-3} \\
    Mat\'ern$(1/2)$      & $\xi(\mu)$  & $0.606^{2}\cdot\mathrm{Mat\acute{e}rn}(\ell=1.60,\nu=1/2)$ & $0.642$ & \num{1.74e-3} \\
    RQ                  & $\xi(\mu)$  & $0.381^{2}\cdot{\rm RQ}(\alpha=10^{5},\ell=2.32)$ & $0.867$ & \num{4.80e-6} \\
    Combined            & $\xi(\mu)$  & $0.406^{2}\cdot\mathrm{Mat\acute{e}rn}(\ell=1160)\cdot{\rm RBF}(\ell=4.86)$ & $0.930$ & \num{3.89e-4} \\
    \bottomrule
  \end{tblr}
\end{table}

\textbf{注意：} （1）RBF 核在 $\xi(\mu)$ 上的 WhiteKernel 噪声 ${}=0.993$（占总方差
$99.93\%$），信号核方差仅为 $0.026^{2}=0.0007$，仿真器完全由噪声主导，无法学习参数依赖性
（详见 \autoref{sec:joint-info}）。（2）$\xi(\mu)$ 上所有核的 WhiteKernel 噪声均偏高
（$0.642$--$0.993$），反映该统计量信号弱、GP 建模困难。（3）Mat\'ern$(1/2)$ 核在
$\xi(s)$ 上的 length\_scale ${}=281$ 远大于其他核（$11$--$60$），说明过粗糙的先验迫使 GP
在极大尺度上平滑，丧失了局部拟合能力。（4）RQ 核的 $\alpha=1.06$（$\xi(s)$），接近 $1$
而远非无穷大，说明其确实使用了多尺度混合特性；但 $\xi(\mu)$ 上 $\alpha=10^5$ 退化为 RBF。

\section{统计显著性检验}\label{sec:significance}

对 LOO-CV 中每个样本的 MAE 进行配对 Wilcoxon\cite{Wilcoxon1945} 符号秩检验，验证核函数
差异的统计显著性。\autoref{tab:wilcoxon-vs-m52} 列出以 Mat\'ern$(\nu=5/2)$ 为基准的对比
结果，\autoref{tab:wilcoxon-pairs} 列出完整配对矩阵中的关键结果。

\begin{table}[H]
  \caption{核函数精度差异：配对 Wilcoxon 检验（vs Mat\'ern$(\nu=5/2)$ 基准）}
  \label{tab:wilcoxon-vs-m52}
  \centering
  \begin{tblr}{lcccc}
    \toprule
    对比                                         & $\xi(s)$ $p$ 值 & 显著性 & $\xi(\mu)$ $p$ 值 & 显著性 \\
    \midrule
    RBF vs Mat\'ern$(5/2)$                        & ${<}0.0001$    & $\ast\ast\ast$ & ${<}0.0001$ & $\ast\ast\ast$ \\
    Mat\'ern$(3/2)$ vs Mat\'ern$(5/2)$             & $0.8114$       & n.s.           & ${<}0.0001$ & $\ast\ast\ast$ \\
    Mat\'ern$(1/2)$ vs Mat\'ern$(5/2)$             & ${<}0.0001$    & $\ast\ast\ast$ & ${<}0.0001$ & $\ast\ast\ast$ \\
    RQ vs Mat\'ern$(5/2)$                          & $0.0003$       & $\ast\ast\ast$ & $1.0000$    & n.s.           \\
    Combined vs Mat\'ern$(5/2)$                    & $0.3955$       & n.s.           & $0.0004$    & $\ast\ast\ast$ \\
    \bottomrule
  \end{tblr}
\end{table}

\begin{table}[H]
  \caption{$\xi(s)$ 上的关键配对比较}
  \label{tab:wilcoxon-pairs}
  \centering
  \begin{tblr}{lcc}
    \toprule
    对比                          & $\xi(s)$ $p$ 值 & 显著性          \\
    \midrule
    RBF vs Mat\'ern$(3/2)$         & $0.0037$        & $\ast\ast$       \\
    RBF vs RQ                     & ${<}0.0001$     & $\ast\ast\ast$   \\
    Mat\'ern$(3/2)$ vs RQ          & $0.0438$        & $\ast$           \\
    Mat\'ern$(1/2)$ vs RQ          & ${<}0.0001$     & $\ast\ast\ast$   \\
    Mat\'ern$(5/2)$ vs Mat\'ern$(3/2)$ & $0.8114$    & n.s.             \\
    \bottomrule
  \end{tblr}
\end{table}

注：$\ast\ast\ast$ $p<0.001$；$\ast\ast$ $p<0.01$；$\ast$ $p<0.05$；n.s.\ 不显著。
双侧检验。经 Bonferroni 校正（$\alpha=0.05/15=0.0033$，共 $C(6,2)=15$ 对比较），
$\xi(s)$ 上所有 ``$\ast\ast\ast$'' 和 ``$\ast\ast$'' 结果在校正后仍然显著。
\textbf{$\xi(\mu)$ 列的显著性结论在无泄露对照实验中基本消失
（见 \autoref{sec:hpo}），仅作为信息泄露影响的定量证据，不作为核函数推荐依据。}

结论：（1）$\xi(s)$ 上形成明确的性能层级：
$\mathrm{Mat\acute{e}rn}(5/2)\approx\mathrm{Mat\acute{e}rn}(3/2)\approx{\rm Combined}
> {\rm RQ} > {\rm RBF} \gg \mathrm{Mat\acute{e}rn}(1/2)$；
（2）Mat\'ern$(5/2)$ 与 Mat\'ern$(3/2)$ 无统计显著差异（$p=0.81$），表明 $\nu=5/2$ 和
$3/2$ 均为合适选择；（3）Mat\'ern$(1/2)$ 显著最差（MAE 比 RBF 高 $45\%$），说明过低光滑度
不适合宇宙学仿真；（4）RQ 核优于 RBF 但不及 Mat\'ern 族，其多尺度混合特性未带来额外优势；
（5）$\xi(\mu)$ 上差异很小且在无泄露条件下不显著，\textbf{本文不将 $\xi(\mu)$ 的核函数
排序作为推荐依据}，核函数选择结论仅基于 $\xi(s)$ 上的 MAE 比较。综合考虑精度和模型
简洁性，推荐使用 Mat\'ern$(\nu=5/2)$ 核。

\section{MAE 的物理含义}\label{sec:mae-meaning}

\autoref{tab:mae-physical} 将 MAE 与数据本身的统计误差对比，判断仿真器精度是否满足科学
分析要求。

\begin{table}[H]
  \caption{MAE 物理含义分析（Mat\'ern 核）}
  \label{tab:mae-physical}
  \centering
  \begin{tblr}{lcc}
    \toprule
    指标                            & $\xi(s)$              & $\xi(\mu)$              \\
    \midrule
    MAE                             & $0.00514$             & $0.01532$               \\
    MAE / $|\mathrm{data}|$         & $0.51\%$              & $1.53\%$                \\
    MAE / $\sigma_{\rm stat}$       & $3.4\%$（${\ll}1$）   & $15.8\%$（偏高）         \\
    MAE / data\_std                 & $7.4\%$               & $77.8\%$                \\
    \bottomrule
  \end{tblr}
\end{table}

$\xi(s)$ 的仿真器误差仅为统计误差的 $3.4\%$，满足精密宇宙学分析的一般要求（${<}5\%$）。
$\xi(\mu)$ 的仿真器误差达统计误差的 $15.8\%$，高于精密宇宙学的理想标准（通常要求 ${<}5\%$）。
更重要的是，MAE/data\_std ${}=77.8\%$ 意味着仿真器误差接近数据在参数空间中的总变化幅度，
部分核函数（尤其 RBF）的 $\xi(\mu)$ 仿真器几乎无法区分不同参数点产生的信号
（详见 \autoref{sec:joint-info} 诊断分析）。因此基于 $\xi(\mu)$ 仿真器得出的参数约束结论
需非常谨慎地解释。主要原因是 $\xi(\mu)$ 信号的参数依赖性极弱
（原始变化范围 $0.87$--$1.28$），仿真器学习这种微弱依赖性存在固有困难。

\section{逐 bin 相对误差分析}\label{sec:per-bin}

\autoref{fig:per-bin} 展示了六种核函数在每个 $s/\mu$ bin 上的平均相对误差。左图 $x$ 轴为
实际物理分离 $s\in[0.5,119.5]~\mathrm{Mpc}/h$（由配置文件 ${\rm s\_range}=[0,120]$ 确定，
120 个 bin）；右图 $x$ 轴为 $\mu\in[0,1]$（97 个 bin）。对于 $\xi(s)$，在大尺度
（$s>100~\mathrm{Mpc}/h$）处相对误差增大，原因是关联函数值接近零导致分母趋小。

\begin{figure}[H]
  \centering
  \includegraphics[width=0.95\linewidth]{image4.png}
  \caption{逐 bin 平均相对误差（左：$\xi(s)$；右：$\xi(\mu)$）}
  \label{fig:per-bin}
\end{figure}

\section{预测对比示例}\label{sec:pred-vs-truth}

\autoref{fig:pred-vs-truth} 展示了六种核函数在两组测试样本上的预测值与 N 体模拟真值的对比。
除 Mat\'ern$(1/2)$ 外，其余核函数均能准确复现关联函数形状。

\begin{figure}[H]
  \centering
  \includegraphics[width=0.95\linewidth]{image5.png}
  \caption{仿真器预测与 N 体真值对比（上行：$s^{2}\xi(s)$；下行：$\xi(\mu)$）}
  \label{fig:pred-vs-truth}
\end{figure}

\section{MCMC 参数约束}\label{sec:mcmc-results}

\textbf{本节结果的三点重要说明}：（1）协方差矩阵由边长 $400~\mathrm{Mpc}/h$ 模拟箱估计，
而训练数据来自 $1000~\mathrm{Mpc}/h$ 箱，体积不匹配导致协方差被高估，
$\chi^2/\mathrm{dof}$ 远小于 $1$，后验分布显著受先验形状驱动（详见
\autoref{sec:omega-c-bias}）；因此本节所有 FoM 的绝对数值不具直接物理意义，仅作核函数
\textbf{相对排序}的依据。（2）MCMC 使用的``观测值''来自 $600~\mathrm{Mpc}/h$ 模拟箱，与
训练数据箱体不同，约 $51.5\%$ 的 $\xi(\mu)$ bin 和 $16.7\%$ 的 $\xi(s)$ bin 落在训练数据
取值范围之外，仿真器在这些 bin 上是外推而非内插。（3）本文仅拟合 $\Omega_c$ 和 $\sigma_8$
两个参数，其余 $7$ 个固定在 fiducial 值；实际 CSST 分析需边缘化所有参数，FoM 将显著降低。
以上三点同时影响所有六种核函数，因此不影响相对排序结论，但读者不应把本节数字当作 CSST
的实际预期约束能力。

\autoref{fig:contour-by-kernel} 和 \autoref{fig:contour-by-mode} 分别按核函数和约束模式
分组展示 MCMC 后验分布。\autoref{tab:mcmc} 列出约束结果汇总。

\begin{figure}[H]
  \centering
  \includegraphics[width=0.95\linewidth]{image6.png}
  \caption{按核函数分组的 $\Omega_c$--$\sigma_8$ 约束（$68\%$ 和 $95\%$ 置信区间）}
  \label{fig:contour-by-kernel}
\end{figure}

\begin{figure}[H]
  \centering
  \includegraphics[width=0.95\linewidth]{image7.png}
  \caption{按约束模式分组的核函数比较}
  \label{fig:contour-by-mode}
\end{figure}

\begin{table}[H]
  \caption{MCMC 参数约束的相对比较（六种核函数；FoM 绝对值不具物理意义，仅作相对排序）}
  \label{tab:mcmc}
  \centering
  \footnotesize
  \begin{tblr}{llccc}
    \toprule
    核函数              & 模式        & $\Omega_c$（中位数 $\pm\sigma$） & $\sigma_8$（中位数 $\pm\sigma$） & FoM $\pm$ bootstrap $\sigma$ \\
    \midrule
    RBF                 & $\xi(s)$    & $0.297\pm0.046$ & $0.775\pm0.074$ & $49.3\pm0.1$ \\
    RBF                 & $\xi(\mu)$  & $0.302\pm0.058$ & $0.856\pm0.184$ & $12.9\pm0.0$ \\
    RBF                 & Joint       & $0.297\pm0.046$ & $0.775\pm0.074$ & $49.3\pm0.1$ \\
    \midrule
    Mat\'ern$(5/2)$      & $\xi(s)$    & $0.300\pm0.045$ & $0.754\pm0.081$ & $47.2\pm0.1$ \\
    Mat\'ern$(5/2)$      & $\xi(\mu)$  & $0.268\pm0.050$ & $0.644\pm0.107$ & $26.1\pm0.1$ \\
    Mat\'ern$(5/2)$      & Joint       & $0.299\pm0.041$ & $0.721\pm0.070$ & $58.9\pm0.2$ \\
    \midrule
    Mat\'ern$(3/2)$      & $\xi(s)$    & $0.294\pm0.047$ & $0.762\pm0.087$ & $43.9\pm0.1$ \\
    Mat\'ern$(3/2)$      & $\xi(\mu)$  & $0.266\pm0.050$ & $0.645\pm0.107$ & $26.1\pm0.1$ \\
    Mat\'ern$(3/2)$      & Joint       & $0.291\pm0.040$ & $0.727\pm0.072$ & $58.3\pm0.2$ \\
    \midrule
    Mat\'ern$(1/2)$      & $\xi(s)$    & $0.292\pm0.047$ & $0.721\pm0.141$ & $32.7\pm0.1$ \\
    Mat\'ern$(1/2)$      & $\xi(\mu)$  & $0.267\pm0.049$ & $0.654\pm0.110$ & $25.8\pm0.1$ \\
    Mat\'ern$(1/2)$      & Joint       & $0.296\pm0.037$ & $0.675\pm0.091$ & $56.3\pm0.2$ \\
    \midrule
    RQ                  & $\xi(s)$    & $0.304\pm0.045$ & $0.752\pm0.081$ & $47.2\pm0.1$ \\
    RQ                  & $\xi(\mu)$  & $0.271\pm0.052$ & $0.639\pm0.103$ & $25.9\pm0.1$ \\
    RQ                  & Joint       & $0.304\pm0.041$ & $0.715\pm0.072$ & $57.7\pm0.2$ \\
    \midrule
    Combined            & $\xi(s)$    & $0.300\pm0.045$ & $0.755\pm0.083$ & $46.6\pm0.1$ \\
    Combined            & $\xi(\mu)$  & $0.277\pm0.056$ & $0.645\pm0.128$ & $19.3\pm0.1$ \\
    Combined            & Joint       & $0.307\pm0.044$ & $0.721\pm0.078$ & $50.4\pm0.1$ \\
    \bottomrule
  \end{tblr}

  \vspace{0.3em}
  {\footnotesize 注：表中 FoM 的 bootstrap 标准差反映 MCMC 链的有限采样误差，不包含
  协方差矩阵高估、箱体不匹配、仅拟合 2 参数等系统性来源的不确定性；因此 FoM 之间
  $<1$ 的差异不宜直接作为核函数显著性判据。核函数推荐主要以 $\xi(s)$ 上的 LOO MAE
  （\autoref{tab:loo}、\autoref{tab:wilcoxon-vs-m52}）为依据。}
\end{table}

\section{MCMC 收敛诊断}\label{sec:convergence}

\textbf{（1）Gelman--Rubin $\hat R$。} 所有 18 组 MCMC 运行的 $\hat R$ 均 ${<}1.01$
（\autoref{tab:rhat}），远低于收敛阈值 $1.1$。

\textbf{（2）自相关时间。} 约 $30$--$47$ 步。有效样本数按 ${\rm N\_eff}=({\rm N\_steps}-
{\rm burnin})/\tau$ 估计（不乘以 walker 数，因 emcee 的 walker 间非独立），约为
$64$--$99$ 个独立样本，对二维参数的后验估计是充分的。\autoref{tab:rhat} 中``有效样本''
按保守估计 ${\rm N\_steps}/\tau_{\max}$ 计算。

\textbf{（3）链轨迹图。} \autoref{fig:trace} 展示六种核函数在 $\xi(s)$ 模式下的 $\Omega_c$
参数 MCMC 链轨迹（3000 步，1000 步 burn-in，配置与正文 5000 步运行相同但步数较短；
红色虚线标记 burn-in 边界），所有核函数的链均在 burn-in 后迅速混合并稳定，无漂移趋势，
验证了收敛性。

\begin{figure}[H]
  \centering
  \includegraphics[width=0.95\linewidth]{image8.png}
  \caption{MCMC 链轨迹图（六种核函数，$\xi(s)$ 模式，$\Omega_c$ 参数，
    40 条 walker $\times$ 3000 步含 1000 步 burn-in；红色虚线为 burn-in 边界）}
  \label{fig:trace}
\end{figure}

\begin{table}[H]
  \caption{MCMC 收敛诊断（Gelman--Rubin $\hat R$，全部六种核函数）}
  \label{tab:rhat}
  \centering
  \footnotesize
  \begin{tblr}{llcccc}
    \toprule
    核函数             & 模式        & $\hat R(\Omega_c)$ & $\hat R(\sigma_8)$ & 自相关时间 & N$_{\rm eff}$ \\
    \midrule
    RBF                 & $\xi(s)$    & $1.004$ & $1.004$ & ${\sim}31$ & ${\sim}97$ \\
    RBF                 & $\xi(\mu)$  & $1.009$ & $1.010$ & ${\sim}47$ & ${\sim}64$ \\
    RBF                 & Joint       & $1.004$ & $1.004$ & ${\sim}31$ & ${\sim}97$ \\
    Mat\'ern$(5/2)$      & $\xi(s)$    & $1.005$ & $1.007$ & ${\sim}31$ & ${\sim}97$ \\
    Mat\'ern$(5/2)$      & $\xi(\mu)$  & $1.008$ & $1.004$ & ${\sim}38$ & ${\sim}79$ \\
    Mat\'ern$(5/2)$      & Joint       & $1.007$ & $1.007$ & ${\sim}32$ & ${\sim}94$ \\
    Mat\'ern$(3/2)$      & $\xi(s)$    & $1.003$ & $1.005$ & ${\sim}32$ & ${\sim}94$ \\
    Mat\'ern$(3/2)$      & $\xi(\mu)$  & $1.004$ & $1.005$ & ${\sim}40$ & ${\sim}75$ \\
    Mat\'ern$(3/2)$      & Joint       & $1.006$ & $1.003$ & ${\sim}33$ & ${\sim}91$ \\
    Mat\'ern$(1/2)$      & $\xi(s)$    & $1.005$ & $1.004$ & ${\sim}33$ & ${\sim}91$ \\
    Mat\'ern$(1/2)$      & $\xi(\mu)$  & $1.004$ & $1.006$ & ${\sim}40$ & ${\sim}76$ \\
    Mat\'ern$(1/2)$      & Joint       & $1.007$ & $1.006$ & ${\sim}32$ & ${\sim}94$ \\
    RQ                  & $\xi(s)$    & $1.004$ & $1.005$ & ${\sim}30$ & ${\sim}99$ \\
    RQ                  & $\xi(\mu)$  & $1.008$ & $1.005$ & ${\sim}40$ & ${\sim}74$ \\
    RQ                  & Joint       & $1.006$ & $1.004$ & ${\sim}32$ & ${\sim}94$ \\
    Combined            & $\xi(s)$    & $1.004$ & $1.007$ & ${\sim}33$ & ${\sim}91$ \\
    Combined            & $\xi(\mu)$  & $1.007$ & $1.010$ & ${\sim}47$ & ${\sim}64$ \\
    Combined            & Joint       & $1.007$ & $1.005$ & ${\sim}31$ & ${\sim}97$ \\
    \bottomrule
  \end{tblr}
\end{table}

\section{重复性验证（3-seed 稳定性）}\label{sec:reproducibility}

为验证超参数优化的可重复性，采用 3 个不同随机种子（${\rm seed}=42$、$123$、$777$）独立
重复训练所有仿真器。\autoref{tab:3seed} 展示完整结果。

\begin{table}[H]
  \caption{3-seed 重复性验证结果}
  \label{tab:3seed}
  \centering
  \footnotesize
  \begin{tblr}{llcc}
    \toprule
    核函数             & 种子            & $\xi(s)$ LOO MAE        & $\xi(\mu)$ LOO MAE       \\
    \midrule
    RBF                 & $42$            & $0.005647$              & $0.015682$               \\
                        & $123$           & $0.005647$              & $0.015300$               \\
                        & $777$           & $0.005647$              & $0.015300$               \\
                        & 均值 $\pm$ 标准差 & $0.005647\pm0.000000$ & $0.015427\pm0.000180$    \\
    \midrule
    Mat\'ern$(5/2)$      & $42$            & $0.005139$              & $0.015320$               \\
                        & $123$           & $0.005139$              & $0.015320$               \\
                        & $777$           & $0.005139$              & $0.015320$               \\
                        & 均值 $\pm$ 标准差 & $0.005139\pm0.000000$ & $0.015320\pm0.000000$    \\
    \midrule
    Mat\'ern$(3/2)$      & $42$            & $0.005088$              & $0.015328$               \\
                        & $123$           & $0.005108$              & $0.015328$               \\
                        & $777$           & $0.005108$              & $0.015328$               \\
                        & 均值 $\pm$ 标准差 & $0.005101\pm0.000010$ & $0.015328\pm0.000000$    \\
    \midrule
    Mat\'ern$(1/2)$      & $42$            & $0.008211$              & $0.015344$               \\
                        & $123$           & $0.008211$              & $0.015344$               \\
                        & $777$           & $0.008211$              & $0.015344$               \\
                        & 均值 $\pm$ 标准差 & $0.008211\pm0.000000$ & $0.015344\pm0.000000$    \\
    \midrule
    RQ                  & $42$            & $0.005392$              & $0.015300$               \\
                        & $123$           & $0.005392$              & $0.015300$               \\
                        & $777$           & $0.005392$              & $0.015410$               \\
                        & 均值 $\pm$ 标准差 & $0.005392\pm0.000000$ & $0.015337\pm0.000052$    \\
    \midrule
    Combined            & $42$            & $0.005156$              & $0.015568$               \\
                        & $123$           & $0.005156$              & $0.015300$               \\
                        & $777$           & $0.005156$              & $0.015300$               \\
                        & 均值 $\pm$ 标准差 & $0.005156\pm0.000000$ & $0.015389\pm0.000126$    \\
    \bottomrule
  \end{tblr}
\end{table}

$\xi(s)$ 的种子间标准差 ${<}0.000001$，结果完全稳定；$\xi(\mu)$ 的标准差约
$0.0001$--$0.0002$（RBF: $0.000180$，Mat\'ern: $0.000000$，Combined: $0.000126$），
表明 20 次 Optuna 试验对于 $\xi(\mu)$ 已基本收敛，核函数间的精度差异不受随机性影响。

\chapter{讨论}\label{chap:discussion}

\section{核函数选择的物理解释}\label{sec:physics-interp}

六种核函数的 $\xi(s)$ 性能排序揭示了清晰的物理规律：
$\mathrm{Mat\acute{e}rn}(5/2)\approx\mathrm{Mat\acute{e}rn}(3/2)\approx{\rm Combined}
> {\rm RQ} > {\rm RBF} \gg \mathrm{Mat\acute{e}rn}(1/2)$。从偏差-方差权衡角度理解：
RBF 核（$\nu=\infty$，无限可微）先验过于刚性，在有限训练样本（104 组）下导致过度平滑
（欠拟合），增大模型偏差。Mat\'ern 核较弱的光滑性假设（$5/2$ 为二阶可微，$3/2$ 为
一阶可微）提供更灵活的拟合能力，在偏差与方差间取得更好平衡。但过度降低光滑度
（$\nu=1/2$，不可微）则走向另一极端——函数空间过大，模型方差急剧增大，MAE（$0.00821$）
比 RBF 还高 $45\%$。这表明宇宙学关联函数对参数的依赖是光滑的，但不需要无限可微，
二阶可微（$\nu=5/2$）或一阶可微（$\nu=3/2$）即为最佳光滑度。

Mat\'ern$(5/2)$ 与 Mat\'ern$(3/2)$ 在 $\xi(s)$ 上无统计显著差异（$p=0.81$，MAE 仅差
$0.00003$），本文推荐 $\nu=5/2$ 基于以下考量：（a）$\nu=5/2$ 是 GP 文献中最广泛采用的
Mat\'ern 变体\cite{Rasmussen2006}，便于与已有工作对比；（b）Mat\'ern$(5/2)$ 核函数二阶
可微，数值稳定性优于一阶可微的 $\nu=3/2$（后者在核矩阵条件数较大时更易出现数值问题）；
（c）联合约束 FoM 略优（$58.9$ vs $58.3$），表明在 MCMC 参数约束中 $\nu=5/2$ 有微弱优势。
但我们强调，在实际应用中 $\nu=3/2$ 同样是合理选择。RQ 核（MAE${}=0.00539$）介于 RBF 和
Mat\'ern 之间，其多尺度混合特性在当前数据上未带来显著优势。组合核（Mat\'ern $\times$ RBF）
在 $\xi(s)$ 上与 Mat\'ern 表现几乎相同（$p=0.40$），进一步检查核参数（\autoref{tab:best-kernels}）
揭示退化行为：在 $\xi(s)$ 上 RBF 因子 length\_scale ${}=9.8\times 10^{4}$（常数），
退化为近似 Mat\'ern 核。

值得注意的是，归一化策略对核函数排序有显著影响。灵敏度分析表明，在不做归一化
（${\rm norm}={\rm False}$）的情况下，三种核的 $\xi(\mu)$ MAE 几乎相同（RBF: $0.01552$，
Mat\'ern: $0.01551$，Combined: $0.01550$），核函数差异消失。而启用归一化后
（${\rm norm}={\rm True}$），Mat\'ern 核改善至 $0.01532$ 而 RBF 核反而劣化至 $0.01568$，
差异变得显著。这说明归一化通过改变 GP 看到的数据尺度，暴露了不同核函数在建模弱参数依赖
性时的能力差异。因此，数据预处理策略本身也是仿真器设计中需要审慎选择的因素。

\section{\texorpdfstring{$\Omega_c$}{Ωc} 系统偏移与协方差矩阵问题}\label{sec:omega-c-bias}

所有 MCMC 约束的 $\Omega_c$ 中位值在 $0.268$--$0.308$ 之间，而 fiducial 值为 $0.26072$，
存在约 $0.01$--$0.05$ 的偏移。后验预测检验（PPC）显示，$\chi^{2}/{\rm dof}$ 在 fiducial
处 $\xi(s)$ 约为 $0.10$、$\xi(\mu)$ 约为 $0.72$（均小于 $1$），说明协方差矩阵被高估
（$\xi(s)$ 严重高估约 $10$ 倍，$\xi(\mu)$ 高估约 $1.4$ 倍），似然面偏平坦，后验分布主要
由先验形状驱动，$\Omega_c$ 的偏移在统计上并不显著。进一步地，各组约束的 $\Omega_c$ 中位值
（$0.297$--$0.308$）极接近先验均值（$(0.2+0.4)/2=0.300$），RBF $\xi(\mu)$ 的 $\sigma_8$
中位值（$0.856$）也接近先验均值（$(0.54+1.18)/2=0.860$），证实在 $\chi^{2}/{\rm dof}$
约 $0.1$ 的条件下，后验分布本质上由先验形状决定，数据对参数的实际约束极弱。

协方差矩阵高估有两方面原因：
\begin{enumerate}
  \item \textbf{体积不匹配。} 协方差从边长 $400~\mathrm{Mpc}/h$ 的模拟箱估计
    （体积 $6.4\times 10^{7}~(\mathrm{Mpc}/h)^{3}$），而仿真器数据来自 $1000~\mathrm{Mpc}/h$
    箱（体积 $10^{9}~(\mathrm{Mpc}/h)^{3}$），有效体积相差约 $15.6$ 倍。若对协方差矩阵
    施加体积重标定（乘以 $(400/1000)^{3}=0.064$），$\chi^{2}/{\rm dof}$ 将从约 $0.10$ 升至
    约 $1.6$（$\xi(s)$）和约 $11.3$（$\xi(\mu)$），前者接近合理值，后者仍偏高（可能源于
    仿真器误差或 bin 间相关性估计偏差）。
  \item \textbf{Hartlap 修正}\cite{Hartlap2007}。 当用 $n_{\rm sim}$ 组模拟估计 $p$ 维
    协方差矩阵时，精度矩阵的无偏估计需乘以
    $f_{\rm Hartlap}=(n_{\rm sim}-p-2)/(n_{\rm sim}-1)$，若模拟数量接近 bin 数量该因子
    将显著偏离 $1$。
\end{enumerate}
以上两点意味着本文的 FoM 数值不具有绝对意义，不代表 CSST 的实际约束能力。由于协方差矩阵
对所有核函数完全相同，核函数间 FoM 的相对排序仍然有效——但需注意：当
$\chi^{2}/{\rm dof}\ll 1$（似然面极平坦）时，FoM 差异可能部分反映仿真器在先验边界附近的
外推行为差异，而非纯粹的核函数预测能力差异。本文的 MAE 比较（\autoref{tab:loo}、
\autoref{tab:wilcoxon-vs-m52}）不受此问题影响，是核函数性能排序最可靠的依据。在实际
CSST 数据分析中，协方差矩阵应从与观测体积匹配的模拟中估计，并施加 Hartlap 修正。

此外，$\sigma_8$ 的偏移方向在不同核函数和约束模式间不一致：使用 $\xi(s)$ 的约束 $\sigma_8$
偏低约 $0.5$--$0.7\sigma$（如 RBF $\xi(s)$: $0.775$，偏离 fiducial $0.810$ 约 $-0.47\sigma$），
而 RBF 的 $\xi(\mu)$ 单独约束偏高（$0.856$，$+0.25\sigma$），Mat\'ern 的 $\xi(\mu)$ 则
大幅偏低（$0.644$，$-1.56\sigma$）。这种方向不一致性进一步表明，当前约束更多反映了
协方差矩阵高估和仿真器外推误差的系统效应，而非真实的参数约束。所有偏移均在 $2\sigma$
以内，在统计上不显著。

\section{联合约束与信息冗余分析}\label{sec:joint-info}

RBF 核的联合约束 FoM（$49.3$）与 $\xi(s)$ 单独完全一致，这并非代码错误，而是因为 RBF 核的
$\xi(\mu)$ 仿真器完全丧失了参数敏感性。诊断分析表明：RBF 核的 $\xi(\mu)$ 模型中
WhiteKernel 噪声项占总方差的 $99.93\%$（${\rm noise\_level}=0.993$），信号核方差仅为
$0.026^{2}=0.0007$，且 RBF length\_scale 被推到极小值（$0.011$），无法建立有效的空间相关性。
结果是仿真器对任意参数输入均给出恒定的预测值（等效于训练集加权平均），$\xi(\mu)$ 的似然
函数在参数空间中为常数，对 MCMC 后验无任何约束贡献。

相比之下，Mat\'ern 核的 $\xi(\mu)$ 仿真器保留了微弱但非零的参数敏感性（$\chi^{2}$ 随
$\Omega_c$ 变化范围约 $2.3$），使其联合 FoM（$58.9$）相比 $\xi(s)$ 单独（$47.2$）提升
$25\%$。这揭示了核函数选择不仅影响预测精度（MAE），更可能决定仿真器能否``看到''参数
依赖性本身。

需要直面的问题是：在当前数据条件下（逐样本均值归一化、仅 129 组训练样本），$\xi(\mu)$ 的
GP 仿真器是否根本不可行？答案是有条件的：Mat\'ern 核证明了弱参数依赖性仍可被学习
（FoM 从 $47.2$ 提升至 $58.9$），但 RBF 核的完全失效表明这极度依赖核函数选择。改善方向
包括：（1）取消逐样本均值归一化以保留振幅信息；（2）采用 ARD 核让 GP 学习参数敏感度差异；
（3）增加训练样本数量。在这些改进实现之前，基于当前 $\xi(\mu)$ 仿真器的 MCMC 约束结论
应谨慎解读。

对原始 129 组 N 体模拟数据（非仿真器预测）的 PCA 分析揭示了信息结构：$\xi(s)$ 的参数
依赖性高度集中在第一主成分（解释 $93.3\%$ 方差），而 $\xi(\mu)$ 更分散（PC1 仅解释
$13.4\%$ 方差）；两者 PC1 的相关系数约 $-0.47$，表明 $\xi(\mu)$ 包含与 $\xi(s)$ 部分
互补的信息。$\xi(\mu)$ 在联合约束中贡献有限的主要原因是其信噪比低
（仿真器 MAE/统计误差 $=15.8\%$），加之逐样本均值归一化削弱了振幅信息，而非两者信息
完全冗余。

\section{与已有仿真器工作的比较}\label{sec:compare-prior}

Cosmic Emulator\cite{Heitmann2009}报告物质功率谱精度约 $1\%$，
EuclidEmulator2\cite{Knabenhans2019}在 $z=0$ 的精度优于 $1\%$。本文 $\xi(s)$ 仿真器 MAE
约 $0.5\%$（信号幅度占比），达到或超过同等精度水平。$\xi(\mu)$ 精度 $1.5\%$ 略低，但
考虑其信号动态范围极小（$0.87$--$1.28$），这是合理的。需指出：不同仿真器的统计量、
参数维度和训练集规模各异，精度数字不宜直接比较。本文贡献在于为给定仿真器框架提供核函数
选择的定量依据。

\section{计算可扩展性}\label{sec:scalability}

GP 回归的核矩阵求逆复杂度为 $O(n^{3})$。本文 129 个训练样本完全可行（单次训练约 $2$ 秒）。
当训练集扩展到 1000+ 样本时，可考虑：（1）稀疏 GP（FITC/SVGP），通过诱导点将复杂度降至
$O(nm^{2})$；（2）随机傅里叶特征（RFF），将核近似为有限维内积；（3）神经网络仿真器
（MLP/CNN），在大数据量下可能更具优势。本文核函数比较结论（Mat\'ern 优于 RBF）在上述
近似方法中的适用性值得进一步研究。

\section{局限性}\label{sec:limitations}

\begin{enumerate}
  \item \textbf{参数约束简化。} 仅拟合 $\Omega_c$ 和 $\sigma_8$ 两个参数，余下 7 个固定在
    fiducial 值；实际分析边缘化所有参数后 FoM 将显著降低。本文 FoM 仅作核函数相对比较基准。
  \item \textbf{重复性。} 3-seed 验证显示 $\xi(s)$ 种子间标准差 ${<}10^{-6}$（完全稳定），
    $\xi(\mu)$ 约 $0.0001$--$0.0002$（波动极小）。20 次 Optuna 试验对 $\xi(s)$ 已充分收敛；
    对 $\xi(\mu)$ 而言可通过增加 ${\rm n\_trials}$ 进一步减小波动。
  \item \textbf{协方差矩阵。} 体积不匹配与 Hartlap 修正问题（见 \autoref{sec:omega-c-bias}）
    影响约束结果的绝对可信度，但不影响核函数间的相对比较结论。
  \item \textbf{观测值外推问题。} MCMC 中使用的``观测值''来自 $600~\mathrm{Mpc}/h$ 模拟箱，
    与训练数据的 $1000~\mathrm{Mpc}/h$ 箱不同，导致约 $51.5\%$ 的 $\xi(\mu)$ bin 落在
    训练数据范围之外，仿真器被迫外推。这是 MCMC 约束结果（特别是 $\Omega_c$ 偏移）的
    重要系统误差来源，但不影响核函数间的相对比较。
  \item \textbf{超参数优化的间接信息泄露。} Optuna 在测试集上优化超参数，而 LOO-CV 使用
    该超参数评估，存在间接泄露（详见 \autoref{sec:hpo} 讨论）。
  \item \textbf{MCMC 先验超出训练边界。} MCMC 先验 $\Omega_c\in[0.2,0.4]$，但训练数据
    $\Omega_c$ 仅覆盖 $[0.187,0.352]$，约 $11$--$14\%$ 的后验样本落在训练范围之外，
    仿真器在参数空间中也存在外推。
  \item \textbf{各向同性 length\_scale。} 本文 GP 使用标量 $\ell$（各向同性核），假设
    9 个参数对统计量的影响``等距''。采用 ARD（Automatic Relevance Determination，每维
    独立 length\_scale）可能显著改善性能，是未来重要改进方向。
  \item \textbf{length\_scale 初始值敏感性。} 本文 Optuna 超参数搜索仅优化
    ${\rm constant\_value}$、${\rm noise\_level}$ 和 $\alpha$ 三个参数，length\_scale
    的初始值固定为 $1.0$，其最终取值依赖 sklearn 内层 L-BFGS-B 优化器的调整。不同核函数
    对 length\_scale 初始值的敏感度可能存在差异——例如 Mat\'ern$(\nu=1/2)$ 的似然面可能
    比 RBF 更加多峰，固定初始值可能导致优化器陷入不同局部极值。本文未对此进行系统性消融
    实验（如多起点随机重启），因此不能完全排除 length\_scale 初始值选择对核函数排名的
    潜在影响。未来工作可将 length\_scale 初始值纳入 Optuna 外层搜索空间，或采用多起点
    策略以提高优化的鲁棒性。
  \item \textbf{未考虑实际系统效应。} 选择效应、光纤碰撞、测光红移误差等均未纳入分析。
\end{enumerate}

\chapter{总结与展望}\label{chap:conclusion}

\section{主要结论}\label{sec:conclusions}

本文基于 129 组高精度 N 体模拟，系统比较了 GP 宇宙学仿真器中六种核函数的性能，主要结论
如下：

\begin{enumerate}
  \item $\xi(s)$ 上核函数性能形成明确层级：
    $\mathrm{Mat\acute{e}rn}(5/2)\approx\mathrm{Mat\acute{e}rn}(3/2)\approx{\rm Combined}
    > {\rm RQ} > {\rm RBF} \gg \mathrm{Mat\acute{e}rn}(1/2)$。Mat\'ern 族（$\nu=5/2$ 和
    $3/2$）显著优于 RBF 核（$p<0.0001$），且二者间无统计显著差异（$p=0.81$），表明
    $\nu=5/2$ 和 $3/2$ 均为合适选择。Mat\'ern$(1/2)$ 因过度粗糙表现最差（MAE 比 RBF 高
    $45\%$）。推荐 Mat\'ern$(\nu=5/2)$ 作为默认核函数。$\xi(\mu)$ 上的核函数差异在无泄露
    条件下不显著，需持保留态度。
  \item 光滑度-性能关系呈非单调``倒 U 形''：过高光滑度（RBF，$\nu=\infty$）导致欠拟合，
    过低光滑度（$\nu=1/2$）导致过拟合，$\nu=5/2$ 和 $3/2$ 取得最优偏差-方差平衡。RQ 核和
    组合核（Mat\'ern $\times$ RBF）均未显著超越单纯 Mat\'ern 核，更复杂的核结构并不必然
    带来优势。
  \item $\xi(s)$ 仿真器精度（MAE/$\sigma_{\rm stat}=3.4\%$）满足精密宇宙学要求；
    $\xi(\mu)$ 仿真器精度（$15.8\%$）有待改进，主要受限于逐样本均值归一化削弱了振幅信息
    和低信噪比。
  \item 联合约束 FoM 排序：Mat\'ern$(5/2)$ ($58.9$) ${}>{}$ Mat\'ern$(3/2)$ ($58.3$)
    ${}>{}$ RQ ($57.7$) ${}>{}$ Mat\'ern$(1/2)$ ($56.3$) ${}>{}$ Combined ($50.4$)
    ${}>{}$ RBF ($49.3$)。Mat\'ern$(5/2)$ 的联合 FoM 相比 $\xi(s)$ 单独提升 $25\%$，
    表明 $\xi(\mu)$ 包含互补信息。FoM 绝对值因协方差矩阵高估不具物理意义，MAE 是核函数
    排序最可靠的依据。
  \item 所有 18 组 MCMC 运行（6 核 $\times$ 3 模式，每组保留 3000 后-burn-in 步）均通过
    Gelman--Rubin 收敛诊断（$\hat R<1.01$），有效独立样本数约 $64$--$99$。
  \item 3-seed 重复性验证和无泄露对照实验共同确认了 $\xi(s)$ 上核函数排序结论的稳健性。
\end{enumerate}

\section{未来展望}\label{sec:future-work}

\begin{enumerate}
  \item \textbf{探索更多核函数及组合方式。} 如 Spectral Mixture 核、周期核、以及核函数的
    加法组合（而非乘法），系统评估更广泛的核函数空间。
  \item \textbf{引入 ARD 核。} 采用各向异性核函数（Automatic Relevance Determination），
    为每个参数维度学习独立的 length\_scale，有望显著改善仿真器对不同参数敏感度的适应性。
  \item \textbf{高维参数约束。} 将 MCMC 拟合扩展到全部 9 个参数，评估边缘化参数后核函数
    选择的影响。
  \item \textbf{改进 $\xi(\mu)$ 仿真器。} 探索 PCA 降维、取消逐样本均值归一化等预处理
    策略，保留 $\sigma_8$ 的振幅信息。
  \item \textbf{应用于 CSST 数据。} 将优化后的仿真器与 CSST 模拟数据结合，为 CSST 科学
    目标提供技术支撑。
  \item \textbf{与神经网络仿真器比较。} 系统评估 GP 仿真器与 MLP/CNN 在精度、不确定性
    估计和计算效率上的优劣。
\end{enumerate}



% !TeX root = ../thesis.tex

\begin{acknowledgements}

本论文的完成离不开许多人的帮助和支持。首先，衷心感谢我的指导教师李霄栋教授，感谢您在选题、方法指导和论文修改方面给予的悉心指导和耐心帮助。

感谢课题组的各位师兄师姐在 N 体模拟数据和代码方面的支持，特别感谢 Zhang 等人提供的高精度仿真数据。感谢各位同学在学习和生活中的陪伴与帮助。

最后，感谢我的家人一直以来的理解和支持，是你们让我能够安心完成学业。

\end{acknowledgements}


\end{document}
